% Part: computability
% Chapter: computability-theory
% Section: coding-computations

\documentclass[../../../include/open-logic-section]{subfiles}

\begin{document}

% SEGMENT OLP-0230-S01

\olfileid{cmp}{thy}{cod}
\olsection{கணக்கீடுகளைக் குறியீடாக்குதல்}

ஒவ்வொரு கணக்கீட்டு மாதிரியிலும் பின்வருவனவற்றைச் செய்ய முடியும்:
\begin{enumerate}
\item கணிக்கத்தக்க சார்புகளின் \emph{வரையறைகளை} முறையாக விவரிக்க.
  எடுத்துக்காட்டாக, டியூரிங் பொறி விவரக்குறிப்புகள், மீள்வரையறைகள்
  அல்லது ஒரு நிரலாக்க மொழியிலுள்ள நிரல்கள் இவ்வரையறைகளைத் தருவதாகக்
  கருதலாம்.
\item ஏதேனும் ஒரு வரையறை தரும் சார்பின், கொடுக்கப்பட்ட உள்ளீட்டிற்கான
  கணக்கீட்டின் முழுப் பதிவை விவரிக்க. எடுத்துக்காட்டாக, ஒரு டியூரிங்
  பொறிக் கணக்கீட்டை, கணக்கீட்டின் ஒவ்வொரு படிக்குமான அமைவுகளின்
  தொடரால் (பொறியின் நிலை, நாடாவின் உள்ளடக்கம்) விவரிக்கலாம்.
\item ஒரு கணக்கீட்டின் சாத்தியப் பதிவு, கொடுக்கப்பட்ட வரையறையுடைய
  கணிக்கத்தக்க சார்பு ஒன்று கொடுக்கப்பட்ட உள்ளீட்டில் எவ்வாறு
  கணிக்கப்படும் என்பதன் பதிவுதானா என்று சோதிக்க (அவ்வுள்ளீட்டில்
  சார்பு வரையறுக்கப்பட்டுள்ளது; அதாவது கணக்கீடு நிற்கிறது).
\item ஒரு கணக்கீட்டின் முழுப் பதிவுக்கான அத்தகைய விளக்கத்திலிருந்து,
  கொடுக்கப்பட்ட உள்ளீட்டிற்கான சார்பின் மதிப்பைப் பிரித்தெடுக்க.
  எடுத்துக்காட்டாக, நிற்கும் ஒரு டியூரிங் பொறிக் கணக்கீட்டின்
  மிகக் கடைசிப் படியிலுள்ள நாடாவின் உள்ளடக்கமே மதிப்பு.
\end{enumerate}

குறியீடாக்கத்தைப் பயன்படுத்தி, அறிவுறுத்தல்களைக் குறியீட்டிலிருந்து
கணிக்கத்தக்க வகையில் மீட்கக்கூடியவாறு, ஒரு கணிக்கத்தக்க சார்பின்
ஒவ்வொரு விளக்கத்திற்கும் எண்ணியல் \emph{சுட்டெண்} ஒன்றை ஒதுக்கலாம்.
அதேபோல், ஒரு கணக்கீட்டின் முழுப் பதிவையும் ஒரு தனி எண்ணால்
குறியீடாக்கலாம். இவ்வாறு கிடைக்கும் எண்கணித உறவு “$s$ ஆனது
சுட்டெண்~$e$ உடைய சார்பின் உள்ளீடு~$x$ க்கான கணக்கீட்டுப் பதிவைக்
குறியீடாக்குகிறது” என்பதும், “குறியீடு~$s$ உடைய கணக்கீட்டுத்
தொடரின் வெளியீடு” என்ற சார்பும் கணிக்கத்தக்கவை; உண்மையில் அவை
முதன்மை மீள்வரையறைக்குரியவை.

இந்த அடிப்படை உண்மை மிகவும் வலிமையானது; தேர்ந்தெடுக்கப்பட்ட
கணக்கீட்டு மாதிரியைச் சாராமல், கணிப்புத்தன்மை குறித்த பல வியப்பூட்டும்
முக்கிய முடிவுகளை நிரூபிக்க இது அனுமதிக்கிறது.

\end{document}
